\section{Introduction}\label{sec:intro}

An Internet access point is shared by devices whose needs differ in kind.
A video call needs a modest rate in every second and is worthless without it, although its picture quality can degrade.
An online game needs very little rate but needs it immediately.
A software update or a mailbox synchronisation needs volume and does not care when it arrives.
The operator of a home broadband service or of public Wi-Fi in a caf\'e, a metro station or an airport also sells different products on the same link: a free tier, a premium subscription, pay-as-you-go access.
And almost nothing is known in advance. Devices join and leave, sometimes dozens at once when a train arrives; demand switches on and off; the usable capacity moves with the backhaul and the radio channel.

When capacity is short, the natural reaction is to protect the real-time traffic and let the delay-tolerant traffic wait.
Doing this well runs into three difficulties at once.
The first is \emph{incentives}. With encrypted traffic the access point often cannot tell what an application is. When it relies on markings set by the device or by the application provider, the class of a flow is in effect declared; when it classifies traffic itself, an application can still shape its traffic to look like a call. We take the extreme case, in which the class is whatever the device says, as the working assumption: a rule that is robust there stays robust when a classifier narrows what can be claimed (Section~\ref{sec:related}).
If declaring ``real-time'' buys priority, every bulk transfer will declare it, and the protection disappears; we show that strict priority and standard utility-maximising schedulers both have this defect.
The second is \emph{fairness under churn}. Guarantees stated as long-run averages are of little use to a device that stays for four minutes.
The third is \emph{business constraints}: premium subscribers were promised a service level, pay-as-you-go traffic is revenue, and a venue that advertises free Wi-Fi must deliver some.

This paper proposes a rule that protects real-time traffic \emph{on credit}.
Real-time flows are served first during a shortage, and each elastic device that is passed over accumulates a deficit that records what it would have received had it claimed the real-time treatment.
The deficit acts in two directions: it raises the priority of the creditor, and it raises the price of further priority for the real-time flows of the same tier.
Elastic users care about volume more than about timing, so a device that stays long enough to be repaid loses little, while a real-time flow gets what it cares about, service now.
The two classes are thus offered different bundles that each prefers to the other's, and, for price-taking devices that stay long enough, telling the truth becomes approximately the best policy without any payment changing hands.

Credit-based sharing is not a new idea; it underlies recent resource managers and ``karma'' economies (Section~\ref{sec:related}).
What we add is a formulation of it inside online constrained optimisation, for service classes that differ in the shape of their utilities and for a population that changes continuously, together with guarantees that hold per device.

\paragraph{Contributions.}
\begin{enumerate}[leftmargin=*,itemsep=2pt]
\item \emph{Model.} We formulate online capacity allocation with unknown arrivals, departures, demands and capacity, three service classes, three commercial tiers with long-run business constraints, and a volume-based notion of \emph{class envy} that captures the temptation to misreport (Section~\ref{sec:model}).
\item \emph{Algorithm.} \repay{} is a drift-plus-penalty rule in which tier-level virtual queues persist across churn and per-device deficit queues are created at arrival. The per-slot problem is made concave by a perspective reformulation and an affine majorant, and is solved exactly by bisection, each step being linear in the number of devices (Section~\ref{sec:algorithm}).
\item \emph{Deterministic envy bound.} On every sample path, including adversarial ones, each deficit is at most $V w_g\eta_k+\rmax_k$. Hence the class envy of every elastic device over every window of its sojourn is bounded by a constant, that does not depend on the number of devices, on the churn, or on the tier prices (Theorem~\ref{thm:bound}).
\item \emph{Incentives.} For price-taking devices the gain in volume from a static misreport is at most this constant, and the gain collected while posing as real-time under a switching strategy is at most the constant times the number of switches; a bound on envy cannot control declarations that change in every slot, which motivates a hold-down on class changes (Theorems~\ref{thm:static}--\ref{thm:dynamic}, Proposition~\ref{prop:commit}). These statements assume price-taking devices and are informative for long stays; we test them against exact deviations in simulation, including in very small markets.
\item \emph{Utility.} For an arbitrary exogenous population process with exogenous demand caps and capacity that is conditionally independent given the population, the time-average utility of the fluid allocation is at most $O(1/V)$ below a benchmark of stationary policies that meet the business constraints in every population state and are class-envy-free in expectation; the violation of the business constraints is $O(\sqrt{V/T})$ in expectation and vanishes almost surely (Theorem~\ref{thm:opt}).
\item \emph{Evaluation.} In simulations of a public hotspot and a home network we compare \repay{} with proportional fairness, weighted fair queueing, strict priority, plain drift-plus-penalty, a Karma-style credit scheduler, strict priority with a capacity cap, and adaptations of six further schedulers and fair-division rules, both when devices are truthful and when they respond strategically, on synthetic traces with varied traffic mixes and population sizes and in three trace-driven scenarios built from public Wi-Fi session logs and mobile throughput logs (Section~\ref{sec:experiments}).
\end{enumerate}

We also state what the approach does not do. It is not the best rule on average performance: in our simulations strict priority with a cap on the real-time share that suits the traffic mix serves about as many real-time slots at a higher elastic rate, and in a trace-driven scenario with little real-time traffic the priority rules serve $12$ to $15$ points more real-time slots; what \repay{} offers instead is a bound for every device that does not depend on the traffic mix. On those traces our default configuration is also not the best one: the variant without the debt price at a smaller $V$ serves more real-time slots with less envy and a similar incentive to misreport. When capacity is short permanently, bounded envy forces real-time flows to be rationed; the benefit of \repay{} lies in transient shortages. The incentive guarantee is approximate: it assumes price-taking devices, it is a constant per stay and therefore weak for short visits, and in our simulations between half and two thirds of the devices still gain a small amount by misreporting, a few of them more than the bound, one in a thin premium tier by a factor of twelve. It also presupposes that demand is measured by the access point, which is natural for downlink traffic only.
Section~\ref{sec:discussion} discusses these limits.
